% !TEX program = pdflatex
% Example: Thesis using memoir class
% Compile with: pdflatex main.tex && bibtex main && pdflatex main.tex && pdflatex main.tex
% Requires: memoir class (included in TeX Live / MiKTeX)

\documentclass[12pt, a4paper, oneside]{memoir}

% Fine-grained page layout
\setlrmarginsandblock{3cm}{3cm}{*}
\setulmarginsandblock{3cm}{3cm}{*}
\checkandfixthelayout

% Chapter style
\chapterstyle{default}

\usepackage{graphicx}
\usepackage{amsmath,amssymb}
\OnehalfSpacing

% Self-contained bibliography
\usepackage{filecontents}
\begin{filecontents*}{references.bib}
@book{goodfellow2016,
  author = {Ian Goodfellow and Yoshua Bengio and Aaron Courville},
  title  = {Deep Learning},
  publisher = {MIT Press},
  year   = {2016}
}
@article{hochreiter1997,
  author = {Sepp Hochreiter and J\"urgen Schmidhuber},
  title  = {Long Short-Term Memory},
  journal = {Neural Computation},
  volume = {9},
  number = {8},
  pages  = {1735--1780},
  year   = {1997}
}
\end{filecontents*}

\begin{document}

\frontmatter
\title{A Comprehensive Study of Neural Networks\\
       for Time Series Prediction}
\author{Author Name}
\date{\today}
\maketitle

\tableofcontents

\mainmatter

\chapter{Introduction}
\section{Background}
Time series prediction is a fundamental problem in
machine learning with applications in finance,
healthcare, and engineering.

\section{Objectives}
This thesis aims to develop improved neural network
architectures for time series prediction.

\chapter{Literature Review}
\section{Traditional Methods}
ARIMA, exponential smoothing, and Kalman filters
have been the standard approaches for decades.

\section{Deep Learning Approaches}
Recent work has explored RNNs, LSTMs~\cite{hochreiter1997}, and Transformers
for time series prediction. A comprehensive overview of deep learning
methods is given in~\cite{goodfellow2016}.

\chapter{Proposed Method}
\section{Architecture}
We propose a hybrid architecture combining temporal
convolutional networks with attention mechanisms.

\chapter{Experiments}
\section{Datasets}
We evaluate on four benchmark datasets.

\section{Results}
Our method achieves the lowest MSE across all datasets.

\chapter{Conclusion}
We presented a novel architecture for time series
prediction that outperforms existing methods.

\backmatter
\bibliographystyle{plain}
\bibliography{references}

\appendix
\chapter{Implementation Details}
\section{Hyperparameters}
The model uses a learning rate of 0.001, batch size
of 32, and 100 training epochs.

\end{document}
